% Generated by scripts/assemble.py; corrections belong in source/corrections.json.
\examhead{15}{北京邮电大学 2026 年数学分析试卷}{20}
% p021-b01; source PDF page 21
\par\Needspace{5\baselineskip}
\blockanchor{p021-b01}
\majorquestion{一}求极限 $\lim_{n\to \infty}\sum_{i = 1}^{n}2^{\frac{i}{n}}\frac{1}{n + \frac{1}{i}}.$\par
% p021-b02; source PDF page 21
\par\Needspace{5\baselineskip}
\blockanchor{p021-b02}
\majorquestion{二}求极限 $\lim_{x\to 0}\frac{\int_0^x e^{\frac{t^2}{2}}\cos t\mathrm{d}t - x}{(e^x - 1)^2(1 - \cos x)\arctan x}.$\par
% p021-b03; source PDF page 21
\par\Needspace{5\baselineskip}
\blockanchor{p021-b03}
\majorquestion{三}求函数 $f(x, y) = \left(y + \frac{x}{3}\right)e^{x + y}$ 的极值.\par
% p021-b04; source PDF page 21
\par\Needspace{5\baselineskip}
\blockanchor{p021-b04}
\majorquestion{四}设 $f(x)$ 在 $x = 0$ 处存在二阶导数, 且 $\lim_{x \to 0} \frac{f(x)}{x} = 0$, 证明: $\sum_{n=1}^{\infty} \left| f\left(\frac{1}{n}\right) \right|$ 收敛.\par
% p021-b05; source PDF page 21
\par\Needspace{5\baselineskip}
\blockanchor{p021-b05}
\majorquestion{五}已知平面区域 $D = \{(x, y) \mid x^2 + y^2 \leq 2y\}$, 计算二重积分 $\iint_{D} (x + y)^2 \mathrm{d}x \mathrm{d}y$.\par
% p021-b06; source PDF page 21
\par\Needspace{5\baselineskip}
\par\Needspace{6\baselineskip}
\blockanchor{p021-b06}
\majorquestion{六}计算第二型曲线积分\par
\nopagebreak[4]
% p021-b07; source PDF page 21
\blockanchor{p021-b07}
\displayformula{I = \int_ {L} x y \mathrm{d} x + (x - y) \mathrm{d} y + x ^ {2} \mathrm{d} z.}
% p021-b08; source PDF page 21
\blockanchor{p021-b08}
其中 L 为螺旋线:  $x = a \cos t, y = a \sin t, z = bt,$  从 t = 0 到  $t = \pi, a, b$  为常数.\par
% p021-b09; source PDF page 21
\par\Needspace{5\baselineskip}
\blockanchor{p021-b09}
\majorquestion{七}求 $I = \iint_{S} (x^{2} + y^{2} + 2) \, \mathrm{d}S$, 其中 $S$ 为球面 $x^{2} + y^{2} + z^{2} = 2ax$, $a > 0$.\par
% p021-b10; source PDF page 21
\par\Needspace{5\baselineskip}
\blockanchor{p021-b10}
\majorquestion{八}设  $a_{n} > 0, n = 1, 2, 3, \cdots$ ，若  $\lim_{n \to \infty} \frac{a_{n+1}}{a_n} = l, l$  为常数，证明： $\lim_{n \to \infty} \sqrt[n]{a_n} = l.$\par
% p021-b11; source PDF page 21
\par\Needspace{5\baselineskip}
\blockanchor{p021-b11}
\majorquestion{九}若函数 $f(x)$ 在 $[1, +\infty)$ 上满足利普希茨条件, 即存在某一个常数 $L > 0$, 使对于任意的 $x_1, x_2 \in [1, +\infty)$, 有 $|f(x_1) - f(x_2)| \leq L|x_1 - x_2|$, 证明: $\frac{f(x)}{x}$ 在 $[1, +\infty)$ 上一致连续.\par
% p021-b12; source PDF page 21
\par\Needspace{5\baselineskip}
\par\Needspace{6\baselineskip}
\blockanchor{p021-b12}
\majorquestion{十}设 $f(x)$ 在 $(0, +\infty)$ 上满足 $f(x) = f(x^2)$, 且 $\lim_{x \to 0^{+}} f(x) = \lim_{x \to +\infty} f(x) = f(1)$, 证明:\par
\nopagebreak[4]
% p021-b13; source PDF page 21
\blockanchor{p021-b13}
\displayformula{f (x) \equiv f (1), x \in (0, + \infty).}
% p021-b14; source PDF page 21
\par\Needspace{5\baselineskip}
\blockanchor{p021-b14}
\majorquestion{十一}设 $f(x)$ 在 $(- \delta, \delta)$ 上有 $n + 1$ 阶导数, 且 $f'(0) = f''(0) = \cdots = f^{(n)}(0) = 0$, $f^{(n+1)}(0) \neq 0$, 由微分中值定理, 对 $|x| < \delta$, 存在 $\theta \in (0,1)$, 使得 $f(x) - f(0) = f'(\theta x)x$, 证明: $\lim_{x \to 0} \theta = \frac{1}{\sqrt[n]{n+1}}$.\par
% p021-b15; source PDF page 21
\par\Needspace{5\baselineskip}
\par\Needspace{6\baselineskip}
\blockanchor{p021-b15}
\majorquestion{十二}设 $f(x)$ 在 $[a, b]$ 上有二阶连续导数且 $f''(x) > 0$, 证明:\par
\nopagebreak[4]
% p021-b16; source PDF page 21
\blockanchor{p021-b16}
\displayformula{f \left(\frac {a + b}{2}\right) \leq \frac {1}{b - a} \int_ {a} ^ {b} f (x) \mathrm{d} x \leq \frac {f (a) + f (b)}{2}.}
% p021-b17; source PDF page 21
\par\Needspace{5\baselineskip}
\blockanchor{p021-b17}
\majorquestion{十三}设二元函数 $f(x, y) = |x - y| \varphi(x, y)$，其中 $\varphi(x, y)$ 在点 (0, 0) 的一个邻域内连续，证明：函数 $f(x, y)$ 在 (0, 0) 处可微的充要条件是 $\varphi(0, 0) = 0$.\par
% p021-b18; source PDF page 21
\par\Needspace{5\baselineskip}
\blockanchor{p021-b18}
\majorquestion{十四}设函数 $f_{n}(x) = x n^{k} e^{-nx}$，其中 $k$ 为常数，讨论 $\{f_{n}(x)\}$ 在 $[0, +\infty)$ 上的一致收敛性.\par
% p021-b19; source PDF page 21
\par\Needspace{5\baselineskip}
\blockanchor{p021-b19}
\majorquestion{十五}设函数 $f(x)$ 在区间 $[a, b]$ 上无界, 证明:\par
% p021-b20; source PDF page 21
\blockanchor{p021-b20}
(1) 存在 $\{x_{n}\} \subset [0,1]$, 使得 $\lim_{n\to \infty}f(x_n) = \infty$.\par
% p021-b21; source PDF page 21
\blockanchor{p021-b21}
(2) 存在 $c \in (a, b)$, 使得对任意的 $\delta > 0$, $f(x)$ 在 $(c - \delta, c + \delta) \cap [a, b]$ 无界.\par
